\section[Repository]{Preparing the repository}
\sectionframe{04}{Preparing the repository}
  {What can the repository tell the agent, so that we do not have to?}
  {\faFolderOpen}

% --------------------------------------------------------------- AGENTS.md
\begin{frame}[fragile]{\texttt{AGENTS.md} for the example}
  \begin{columns}[T,onlytextwidth]
    \begin{column}{0.58\textwidth}
\begin{codefile}[basicstyle=\ttfamily\fontsize{7.2}{8.75}\selectfont]{AGENTS.md}
# sem-seg
## Setup and checks
- Environment: conda env `sem-seg`, activated beforehand
- Fast check (CPU, synthetic data): `pytest -q`
- Lint, types: `ruff check`, `mypy src` §\textcolor{udsred}{\sffamily\scriptsize\faArrowLeft\ after every edit}§
## Conventions
- Terms: `GLOSSARY.md` (specimen, micrograph, patch) §\smash{\textcolor{udsblue}{\faLink}}§
- Decisions: `docs/adr/` §\smash{\textcolor{udsblue}{\faLink}}§
- Lengths in micrometres; pixel size from metadata
## Boundaries
- Never write to `data/raw/`; use `data/derived/`
- No cluster or GPU jobs (`sbatch`, `srun`)
- Ask before adding a dependency
## When you finish
- Report changed files, checks run, open assumptions
\end{codefile}
    \end{column}
    \begin{column}{0.39\textwidth}
      % one bracket per section of the file, level with its lines; y = 0 is the top of the file window
      \begin{tikzpicture}[x=1mm,y=1mm,
          what/.style={anchor=west,inner sep=0pt,font=\footnotesize,text=udsdark,align=left,text width=46mm}]
        \useasboundingbox (0,0) rectangle (56,-56);
        \def\first{6.3}\def\pitch{3.075}% the first line's top below the window's top, and the line pitch
        \foreach \a/\b/\t in {2/5/{commands it can run},
            6/9/{what the code cannot show\\{\scriptsize\color{udsblue}\faLink\ pointers to other files}},
            10/13/{requests; section 8 enforces them},14/15/{the report you get}} {
          \pgfmathsetmacro\ya{-\first-(\a-1)*\pitch-0.5}
          \pgfmathsetmacro\yb{-\first-\b*\pitch+0.5}
          \draw[udsgreymid,line width=0.6pt] (0,\ya) -- (1.4,\ya) -- (1.4,\yb) -- (0,\yb);
          \node[what] at (3.6,{(\ya+\yb)/2}) {\t};
        }
      \end{tikzpicture}
    \end{column}
  \end{columns}
  \note{\textbf{16:40--17:50}\par
    AGENTS.md is the one file that tools supporting it load at the start of every session, so it
      holds commands, boundaries and pointers to other files.\par
    Fifteen lines, each preventing a specific mistake: a guessed environment, a wrong test
      command, a decision undone, assumed pixel sizes, writes into raw data, a GPU job or a new
      dependency nobody asked for.\par
    The environment is activated before the agent starts, as section 3 advised, so the commands
      stay short.\par
    The checks set up in section 3 go in here, so the agent runs them after every edit.\par
    Boundaries written here are requests; section 8 enforces them.\par
    Keep it short and specific, and review it like code.}
\end{frame}

% ------------------------------------------------------------ tree example
\begin{frame}[source={OpenAI (2026), Custom instructions with AGENTS.md}]
  {Which instruction files load?}{Codex started in \texttt{src/sem\_seg/data/}}
  \begin{foldertree}[badge at=50,note at=75,row height=3.9]
    \tikzset{loaded/.style={fill=udsblue,text=white},notl/.style={draw=udsgreymid,text=udsgrey}}
    \tfile{0}{\textasciitilde/.codex/AGENTS.md}{global defaults}\tbadge{loaded}{LOADED 1}
    \tdir{0}{sem-seg/}{contains \texttt{.git}}\tbadge{draw=udsdark,text=udsdark}{ROOT}
    \tfile{1}{AGENTS.md}{repository rules}\tbadge{loaded}{LOADED 2}
    \tfile{1}{GLOSSARY.md, docs/adr/}{read when \texttt{AGENTS.md} points there}%
      \tbadge{draw=udsdark,dashed,text=udsdark}{ON DEMAND}
    \tdir{1}{src/sem\_seg/}{}
    \tfile{2}{AGENTS.md}{package rules}\tbadge{loaded}{LOADED 3}
    \tdir{2}{data/}{\texttt{codex --cd src/sem\_seg/data}}\tbadge{fill=udsred,text=white}{WORKING DIR}
    \tfile{3}{AGENTS.md}{replaced by the override}\tbadge{draw=udsred,text=udsred}{IGNORED}\tstrike
    \tfile{3}{AGENTS.override.md}{last, so it wins conflicts}\tbadge{loaded}{LOADED 4}
    \tfile{3}{augment/AGENTS.md}{below the working directory}\tbadge{notl}{NOT LOADED}
    \tfile{2}{models/AGENTS.md}{beside the working directory}\tbadge{notl}{NOT LOADED}
    \tfile{1}{tests/AGENTS.md}{not on the path to the working directory}\tbadge{notl}{NOT LOADED}
  \end{foldertree}
  \vfill
  {\footnotesize Numbers give the prompt order.\par}
  \note{\textbf{17:50--19:10}\par
    How Codex assembles its instructions, on our repository.\par
    We start Codex in the data package.\par
    The global file loads first, then the root AGENTS.md, the package file and the data folder,
      where the override replaces the AGENTS.md beside it.\par
    Later files come later in the prompt and win conflicts.\par
    Started from the root, only files 1 and 2 would load.\par
    Files below or beside the working directory are not loaded at start, although the agent may
      still open them.\par
    The 32 KiB default limit counts instruction text; it is not the model's context window.\par
    Codex builds the chain once per run, so restart it after editing.\par
    Other tools use other names and rules; the appendix has the Codex details.}
\end{frame}

% ---------------------------------------------------- glossary, standards
\begin{frame}[fragile,source={Pocock (2026), mattpocock/skills v1.3.1: domain-modeling, code-review, retro}]{\texttt{GLOSSARY.md} and \texttt{CODING\_STANDARDS.md}}
  \begin{columns}[T,onlytextwidth]
    \begin{column}{0.485\textwidth}
\begin{codefile}[basicstyle=\ttfamily\fontsize{7}{8.5}\selectfont]{GLOSSARY.md\enspace\textcolor{udsgreymid}{\sffamily\scriptsize read when pointed to}}
# Glossary
Specimen: one physical sample (`specimen_id`).
  Avoid: §\textcolor{udsred}{\bfseries sample}§, coupon. §\textcolor{udsred}{\sffamily\scriptsize\faArrowLeft\ the word in our request}§
Area: one imaged region of a specimen.
Micrograph: one SEM image of an area
  (`micrograph_id`). Avoid: image, frame.
Patch: a crop of a micrograph used in
  training (`patch_id`). Avoid: tile.
Fold: specimens held out together; every
  specimen belongs to exactly one fold.
\end{codefile}
    \end{column}
    \begin{column}{0.485\textwidth}
\begin{codefile}[basicstyle=\ttfamily\fontsize{7}{8.5}\selectfont]{CODING\_STANDARDS.md\enspace\textcolor{udsgreymid}{\sffamily\scriptsize read in review}}
# Coding standards
- Every random step takes a `seed`; the
  same seed gives the same result.
- Public functions check their inputs and
  raise `ValueError` naming the problem.
- Tests call public functions on small,
  hand-written fixtures, never real data.
- Expected values come from the fixture,
  never recomputed the way the code does.
- A new way to split the data needs an ADR.
\end{codefile}
    \end{column}
  \end{columns}
  \note{\textbf{19:10--20:20}\par
    Two files the agent reads only when something points to them.\par
    The glossary decides what words mean: our request said sample, and here sample is a word to
      avoid.\par
    The clarifying step of section 6 writes terms into it as soon as they are agreed.\par
    CODING\_STANDARDS.md says how code is written here: seeds, input checks, how tests are
      built; a fixture is the small, hand-made input a test runs on.\par
    In the workflow of section 6 the review step reads it, not the implementing agent, which
      already carries the most context.\par
    Rules that a linter can check belong in CI instead.}
\end{frame}
